\documentclass[11pt]{article}

% ============================================================
%                     CONFIGURACIÓ BÀSICA
% ============================================================

\usepackage[utf8]{inputenc}
\usepackage[a4paper, margin=1.2cm]{geometry}

% ============================================================
%                    MATEMÀTIQUES I FÍSICA
% ============================================================

\usepackage{amsmath}
\usepackage{amssymb}
\usepackage{amsthm}
\usepackage{mathtools}
\usepackage{physics}

% ============================================================
%                    COLORS I CAIXES
% ============================================================

\usepackage[dvipsnames, svgnames, x11names]{xcolor}
\usepackage{tcolorbox}

% ============================================================
%                       CONFIGURACIÓ
% ============================================================

\graphicspath{{./images/}}
\usepackage{hyperref}
\linespread{1.25}

% ============================================================
%                    COMANDAMENTS PROPIS
% ============================================================

\newcommand{\R}{\mathbb{R}}
\newcommand{\Z}{\mathbb{Z}}

% ============================================================
%                         DOCUMENT
% ============================================================

\title{\textbf{Esquema d'Àlgebra}}
\author{\textit{Gerard Cosp Lores}}
\date{}

\begin{document}

\maketitle

\begin{abstract}
Aquest document recull les "Maneres de fer" \textbf{l'assignatura d'Àlgebra}. \footnote{Aquest document s'anirà actualitzant a mesura que avançi el semmestre.} 
\end{abstract}

\begin{tcolorbox}[
    colback=white,
    colframe=black,
    boxrule=0.8pt,
    arc=0pt
]

\textbf{COMPROVACIÓ D'ESPAI I SUBESPAI}

\end{tcolorbox}

\begin{tcolorbox}[
    colback=white,
    colframe=black,
    boxrule=0.8pt,
    arc=0pt
]

\textbf{TROBAR BASE DEL SUBESPAI}

\textbf{BASE DE LA SUMA I L'INTERSECCIÓ}

\end{tcolorbox}

\begin{tcolorbox}[
    colback=white,
    colframe=black,
    boxrule=0.8pt,
    arc=0pt
]

\textbf{TROBAR MATRIU DE L'APLICACIÓ}

\end{tcolorbox}

\begin{tcolorbox}[
    colback=white,
    colframe=black,
    boxrule=0.8pt,
    arc=0pt
]

\textbf{CANVI DE BASES}

\end{tcolorbox}

\end{document}